\section{Conclusion}
\label{sec:conclusion}

We presented GAME-Mal, a malware family classifier that combines
reflection-aware dynamic-trace parsing with a sigmoid-gated
multi-head attention transformer. The gated architecture --- G1
position of Qiu et al.\ --- reaches macro F1 of $0.884$ on the
eight-family UMD subset, $15.9$ percentage points above the
best-swept reproduction of the Markov associative-rule classifier
of D'Angelo et al.\ ($0.725$), while exposing an intrinsic,
forward-pass-native explanation of its decisions through its gate
activations. Fusob, the minority ransomware family, is a useful
case in point: despite contributing only $1.8\%$ of the corpus,
it achieves perfect recall under GAME-Mal whereas the statistical
baseline collapses on it for want of rule-mining evidence.

The full benchmark picture is more nuanced than the headline gap
suggests. A 2-layer BiLSTM trained under the same protocol
reaches macro F1 of $0.894 \pm 0.003$ --- narrowly ahead of
GAME-Mal and statistically tied with a Random Forest over the
same Markov rule features ($0.893 \pm 0.010$). The gated
transformer does not lead on raw accuracy at this corpus scale.
Our matched-prep ablation confirms the gate itself does not
add accuracy: a plain (non-gated) transformer of identical
capacity reaches F1 $0.897$, marginally above GAME-Mal's
$0.883$.

The broader contribution is therefore \emph{methodological rather
than performance-based}. What the gate gives is an always-on
ranking signal over APIs, available in a single forward pass
with no back-propagation, no reference sample, and no additional
compute. The deletion-test faithfulness analysis
(Section~\ref{sec:deletion}) qualifies this claim: the gate
shows clear signal at $k = 20$ (gate-mask degrades true-class
probability by $1.04$ pp more than random-mask) and for families
with moderate confidence (DroidKungFu: $+4.77$ pp at $k = 20$),
but the effect is negligible at $k \in \{5, 10\}$ in aggregate,
partly because Fusob and Jisut achieve baseline probabilities
of $> 0.9999$ --- a ceiling effect that makes any deletion test
uninformative. We therefore describe the gate as providing
\emph{partial faithfulness}: a useful, cheap per-family ranking
of which tokens the attention layer emphasised, without a
strong claim that every individual top-ranked token is
independently causally responsible for the prediction.

We view the core trade as: $\approx 1$ pp macro-F1 in exchange
for inference-time explainability. On security-critical pipelines
where a human analyst must validate every flagged sample, that
trade is often the one worth making. Reporting it honestly ---
including the cases where the gate does not add accuracy and
where the faithfulness validation is partial --- is what lets
future work build on these findings without inheriting inflated
claims.
